\documentclass[11pt]{article}
\usepackage[margin=1in]{geometry}
\usepackage{booktabs}
\usepackage[round]{natbib}
\bibliographystyle{plainnat}

\title{Wage Growth and Broadband Access in Mid-Sized Texas Metros:\\ A Brief Assessment}
\author{Policy Research Division}
\date{July 2026}

\begin{document}
\maketitle

\begin{abstract}
This brief examines how wage growth in mid-sized Texas metropolitan areas
relates to broadband availability and workforce composition. Drawing on
recent county-level evidence \citep{delgado2023}, we summarize wage and
connectivity indicators for a set of illustrative metros and discuss
implications for state infrastructure policy.
\end{abstract}

\section{Introduction}
\label{sec:intro}

Texas's economic story over the past decade has been written largely in its
biggest metros, but the state's mid-sized cities have quietly diverged from
one another in ways that matter for policy. \citet{delgado2023} document
substantial county-level wage divergence within the Texas Triangle, while
\citet{okafor2024} find that rural broadband buildout shifted labor-market
outcomes in the Panhandle. This brief extends that discussion to a set of
illustrative mid-sized metros, asking whether connectivity gaps track the
wage gaps that policymakers most often cite.

Our approach is descriptive rather than causal. We assemble indicators on
wages and broadband subscription for several metros and compare them against
the demographic context described at length by \citet{hemphill2021}. The
headline pattern, summarized in Table~\ref{tab:wages}, is that connectivity
and wage growth move together, though not uniformly.

\section{Wage and Connectivity Indicators}
\label{sec:indicators}

Table~\ref{tab:wages} reports the core indicators. The metros shown are
illustrative and the figures are synthetic, constructed for demonstration
purposes only. Note that the sourcing for each row is documented directly in
the table, and the methodological caveats are discussed in
Section~\ref{sec:methods}.

\begin{table}[ht]
\centering
\caption{Wage growth and broadband subscription, selected metros}
\label{tab:wages}
\begin{tabular}{lrrl}
\toprule
Metro area & Wage growth (\%) & Broadband (\%) & Source \\
\midrule
Nacogdoches      & 47.3 & 61.9 & \citep{delgado2023} \\
Harlingen        & 38.6 & 58.2 & \citep{okafor2024} \\
San Angelo       & 29.4 & 71.7 & \citep{hemphill2021} \\
Texarkana        & 33.8 & 64.1 & \citep{delgado2023} \\
\bottomrule
\end{tabular}
\end{table}

A second question is how these indicators translate into fiscal capacity.
Table~\ref{tab:fiscal} presents synthetic per-capita revenue figures for the
same metros, with a column pointing back to the methods discussion for the
adjustment applied to each row.

\begin{table}[ht]
\centering
\caption{Synthetic per-capita local revenue, selected metros}
\label{tab:fiscal}
\begin{tabular}{lrrl}
\toprule
Metro area & Revenue (\$) & Index & Notes \\
\midrule
Muleshoe        & 8214 & 103.7 & Adjusted per Section~\ref{sec:methods} \\
Brenham         & 6598 & 91.2  & Unadjusted \\
Refugio         & 7443 & 97.6  & Adjusted per Section~\ref{sec:methods} \\
\bottomrule
\end{tabular}
\end{table}

\section{Methods and Caveats}
\label{sec:methods}

All values in the tables above are fabricated for software testing and
should not be quoted. The ``adjustment'' referenced in
Table~\ref{tab:fiscal} is a simple cost-of-living rescaling in the spirit of
the deflators discussed by \citet{hemphill2021}. Broadband subscription
shares follow the definitional conventions used by \citet{okafor2024}, and
the wage-growth window matches \citet{delgado2023}.

\section{Conclusion}
\label{sec:conclusion}

The descriptive evidence in Tables~\ref{tab:wages} and \ref{tab:fiscal}
suggests that connectivity and fiscal capacity vary together across
mid-sized Texas metros, consistent with the broader literature
\citep{delgado2023,okafor2024,hemphill2021}. Causal claims would require a
research design well beyond this brief, but the correlations alone justify
continued attention to broadband policy in the state's middle tier of
cities.

\bibliography{refs}

\end{document}
